\section{Exact source map and reproducibility}\label{sec:sources}
All imports are public text at immutable commits. F and V are in
\texttt{mxym/math} at
\begin{center}\small\texttt{\srcpin}.\end{center}
U is in \texttt{OpenAI/math} at
\begin{center}\small\texttt{\uppin}.\end{center}
The original shorthand \texttt{mxym/math009} did not identify an accessible
repository; the audited source is entry 009 within \texttt{mxym/math}.
The twenty-five audited F/V files have identical Git blob hashes at the
first commit and \path{1c67b4146f9e54291c8a5ca1f42a028bf307629f}.
No result here depends on an unpinned working tree.

\subsection{Imported results and their exact locations}
In the following map, paths are repository-relative and labels are the
literal TeX labels in the pinned source. Line numbers are conveniences for
the pinned bytes, not the basis of identification.

\paragraph{Virial supplement V.}
All of the following labels are in
\path{notes/full-density-virial-stress/manuscript.tex}:
\begingroup\small
\begin{longtable}{@{}p{0.24\textwidth}p{0.40\textwidth}p{0.29\textwidth}@{}}
\toprule
Use here & Exact source label(s) & Pinned line(s) \\
\midrule\endhead
Model, diameter, ensemble & \lab{ass:model}, \lab{eq:ensemble} & 46, 56 \\
Collision and sphere normalization & \lab{eq:collision}, \lab{eq:J} & 70, 82 \\
H1 & \lab{imp:initial} & 142 \\
H2 & \lab{imp:algebra}, \lab{eq:partition}, \lab{eq:rootoutput} & 145, 149, 151 \\
H3 & \lab{imp:unweighted}, \lab{eq:unweighted} & 163, 165 \\
H4 & \lab{imp:majorant}, \lab{eq:majorant} & 172, 179 \\
H5 & \lab{imp:tree} & 184 \\
H6 & \lab{imp:coupling} & 193 \\
Mechanical identities & \lab{lem:mechanics}, \lab{eq:micro-virial}, \lab{eq:micro-spread} & 210, 213, 215 \\
Vacancy expansion & \lab{lem:initial}, \lab{eq:vacancy} & 233, 235 \\
Incidence mean and clipping & \lab{prop:incidence}, \lab{eq:mark-bound}, \lab{eq:clip}, \lab{eq:rawbound}, \lab{eq:scalar-moment} & 257, 286, 296, 308, 323 \\
First-moment transfer & \lab{lem:transfer}, \lab{eq:transfer}, \lab{eq:true-incidence} & 344, 346, 369 \\
\bottomrule
\end{longtable}
\endgroup
The proof of Lemma~\ref{lem:incidence-mean} reproduces the weighted
mean argument in V. V's $H_\eps^p$ denotes an expectation; this manuscript's
$\cH_\eps^p$ is random and $m_\eps=\E\cH_\eps^p$. Confusing these two
would bypass the required variance estimate.

\paragraph{Functional manuscript F.}
Use the prefix
\path{preprints/009-functional-hard-sphere-fluctuations/v1/sections/}.
\begin{longtable}{@{}p{0.31\textwidth}p{0.35\textwidth}p{0.27\textwidth}@{}}
\toprule
File & Exact labels & Use here \\
\midrule\endhead
\path{01_model.tex} & \lab{ass:kinetic}, \lab{lem:convention} & Kinetic interval, path versions \\
\path{02_imports.tex} & \lab{eq:counts}, \lab{prop:imports} & Factorial/cumulant normalization and package \\
\path{05_increments.tex} & \lab{lem:remainder} & Unweighted record errors \\
\path{06_microscopic.tex} & \lab{lem:record} & Passive bounded records and zero continuation \\
\path{07_functional.tex} & \lab{thm:functional} & Exactly centered Schwartz functional limit \\
\path{08_tempered.tex} & \lab{thm:tempered}, \lab{eq:tempered-covariance} & Strong-dual path limit, Gaussian covariance \\
\bottomrule
\end{longtable}
The connected scaling at several roots in \eqref{eq:scaling} uses F's
normalized factorial convention and V's chronological birth/link count.
The initial Gaussian covariance and initial/noise independence are those
of F. Its positive-time test space is retained exactly as stated.

\paragraph{Upstream source U.}
Use the prefix
\path{preprints/Hard-sphere-fluctuations-on-the-regular-Boltzmann-lifespan-September-23-2026/build/sections/}.
\begin{longtable}{@{}p{0.32\textwidth}p{0.35\textwidth}p{0.26\textwidth}@{}}
\toprule
File & Exact labels or portion & Use here \\
\midrule\endhead
\path{02-initial-target.tex} & \lab{prop:initial}, proof's static cumulant bound & Bounded initial CLT and $\mu\|c_{2,0}\|_{\TV}=O(\eps)$ \\
\path{03-dhm-inputs.tex} & Gaussian input envelope and dynamics & Bol inputs within H1 \\
\path{04-cluster-expansion.tex} & Autonomous partition and coarse variation argument & Operational basis of H2/H3 \\
\path{05-cutoff.tex} & \lab{prop:cutoff}, \lab{eq:cutoff-probability} & Full-time $O(\eps^{21/20})$ probability comparison \\
\path{06-sharp-asymptotics.tex} & \lab{eq:sharp-input-bounds}, \lab{eq:sharp-scaling}, \lab{lem:sharp-factorials}, \lab{lem:sharp-summability}, \lab{eq:sharp-energy-sum} & Chronological scaling, energy bins, Stirling and entropy sum \\
\path{07-gaussian-limit.tex} & Leading-tree cancellation and Gaussian identification & Basis of H5 and input F \\
\bottomrule
\end{longtable}
The seven-file U inventory includes its \path{build/main.tex}. These
operational and functional inputs are the conditional foundation; this
manuscript does not claim to reprove the entire upstream hierarchy.

\subsection{Proof contribution and conclusion ledger}
Lemma~\ref{lem:weighted} proves the previously missing degree-two,
two-root passive-speed estimate by global clipping and retained Gaussian
energy. The exact variance formula then proves \eqref{eq:Hvar}; monotonicity
proves the uniform pasted $L^2$ LLN and H6 gives the true probability LLN.
Sections~\ref{sec:mechanics} and \ref{sec:process} supply maximal quadratic
closure, the full thirteen-component theorem, the joint-field argument,
and pasted covariance convergence. Corollary~\ref{cor:intrinsic} adds the
per-particle intrinsic observables via explicit differentials.
Lemma~\ref{lem:endpoint} is a separate endpoint covariance result.
The true prelimit covariance obstruction in
\eqref{eq:failed-L2-transfer} remains an unresolved transfer problem.

\subsection{Reproduction package}
The source archive contains this manuscript and all its sections, the PDF,
\path{build.py}, \path{README.txt}, optimization-safe symbolic checks,
negative-control evidence, and source provenance. The thirty-two exact pinned
verification-cache files are bundled under \path{sources/}; their URL, byte
count, SHA-256 and Git blob SHA-1 are recorded in
\path{provenance/source-manifest.json} and
\path{provenance/upstream-manifest.json}. This cache contains the audited
imports, rather than complete standalone distributions of F, V and U.
The public strict checker verifies all cache bytes and literal source labels
offline; optional fetching uses ordinary TLS verification.

\path{audit/AUDIT.txt} is a documented public derivative of the original
mathematical audit. The seven original audit objects remain unchanged and
are represented by their hashes in \path{provenance/original-object-pins.json}.
The original patch has SHA-256
\begin{center}\small
\path{6c659d81faafc9a803efe97da399e16aa7aa642ab352213e5c3f77d360442528}.
\end{center}
The public checker uses explicit exceptions and verifies fixed audit,
provenance, manuscript and cached-source anchors under both ordinary Python
and \texttt{python -O}. Deliberate corruption is rejected in both modes.
These checks verify identities, normalizations and source correspondence;
the analytic proofs are in the text above. No Lean verification is claimed.

The build fixes the source epoch, omits volatile PDF dates and identifiers,
disables shell escape, and performs three pdfLaTeX passes in a fresh
directory. \path{build/determinism.json} records two fresh matching builds.
The isolated fallback format loads \texttt{english hyphen.tex} when the
standard installation lacks its generated format or font map.
Byte reproducibility across different TeX or font versions is not asserted.
\path{PUBLIC_FILES.json} is the publication whitelist;
\path{package-manifest.json} hashes its files except the self-excluded
manifest. \path{checks/verify_package.py} also rejects unexpected files.
